\ProvidesFile{StringLiteral.tex}[Строковые литералы]


\subsection{Строковые литералы}

Строковые литералы являются заменой голому \verb"const char*" при работе со строками.
Класс позволяет использовать строковые литералы как аргументы шаблона напрямую.
Давайте посмотрим что можно еще выжать из этого класса.

\subsubsection{Базовая имплементация}

Давайте стартуем с имплементации, которая была упомянута в самом начале раздела.
\begin{cppcode}
template<size_t N>
struct StringLiteral {
  constexpr StringLiteral(const char (&str)[N]) {
    std::copy_n(str, N, data);
  }

  char data[N]{};
};
\end{cppcode}
Чтобы класс был используемым во время компиляции, у него все данные должны быть публичные.
Теперь давайте добавим в этот класс некоторые методы, которые обеспечат удобство взаимодействия со строками.
Начнем с методов доступа:
\begin{cppcode}
static constexpr std::size_t capacity() noexcept {
  return N;
}

static constexpr std::size_t size() noexcept {
  return N - 1;
}

constexpr std::string_view view() const {
  return {data, N};
}
\end{cppcode}
Так как в \verb"StringLiteral" строка оканчивается \verb"'\0'", то \verb"N" будет на $1$ больше, чем длина строки.
Если мы захотим произвести какую-то посимвольную модификацию строки, то нам полезно иметь метод
\begin{cppcode}
constexpr StringLiteral map(auto f) const noexcept {
  char tmp[N];
  if constexpr (N > 1)
    std::transform(data, data + N - 1, tmp, f);
  tmp[N - 1] = '\0';
  return tmp;
}
\end{cppcode}
И тогда можно сделать что-то такое
\begin{cppcode}
StringLiteral str1 = "Some Phrase!";
StringLiteral str2 = str1.map(to_lower_case);
// str2 equals "some phrase!"
\end{cppcode}
Здесь подразумевается, что функция \verb"to_lower_case" определена где-то выше и переводит все буквы английского алфавита в строчный формат, а остальные символы не трогает.

Так же полезно будет получать доступ к началу и концу строки, к первому и последнему элементу.
Для этого определим функции \verb"front" и \verb"back", которые отщепляют от строки ее часть сначала или с конца.
Если количество символов, что надо отщепить больше длины строки, то отщипываем всю строку.
\begin{cppcode}
constexpr char front() const noexcept {
  return data[0];
}

template<size_t K>
constexpr StringLiteral<(K < N ? K : N - 1) + 1> front() const noexcept {
  static constexpr size_t K0 = (K < N ? K : N - 1);
  char tmp[K0 + 1];
  std::copy_n(data, K0, tmp);
  tmp[K0] = '\0';
  return tmp;
}
\end{cppcode}
И аналогично функция для откусывания конца строки:
\begin{cppcode}
constexpr char back() const noexcept {
  return data[N - 2];
}

template<size_t K>
constexpr StringLiteral<(K < N ? K : N - 1) + 1> back() const noexcept {
  static constexpr size_t K0 = (K < N ? K : N - 1);
  char tmp[K0 + 1];
  std::copy_n(data + (N - K0 - 1), K0, tmp);
  tmp[K0] = '\0';
  return tmp;
}
\end{cppcode}
Так же полезно иметь функции, которые наоборот указывают сколько символов надо выкинуть сначала или с конца.
Вот как это выглядит для выкидывания элементов в начале:
\begin{cppcode}
template<size_t K>
constexpr StringLiteral<(K < N ? N - K : 1)> pop_front() const noexcept {
  static constexpr size_t K0 = (K < N ? K : N - 1);
  char tmp[N - K0];
  std::copy_n(data + K0, N - K0, tmp);
  return tmp;
}
  
constexpr StringLiteral<N - 1> pop_front() const noexcept {
  return pop_front<1>();
}
\end{cppcode}
А вот для выкидывания элементов в конце:
\begin{cppcode}
template<size_t K>
constexpr StringLiteral<(K < N ? N - K : 1)> pop_back() const noexcept {
  static constexpr size_t K0 = (K < N ? K : N - 1);
  char tmp[N - K0];
  if constexpr (N - K0 > 1)
    std::copy_n(data, N - K0 - 1, tmp);
  tmp[N - K0 - 1] = '\0';
  return tmp;
}
  
constexpr StringLiteral<N - 1> pop_back() const noexcept {
  return pop_back<1>();
}
\end{cppcode}
Теперь можно сделать что-то такое
\begin{cppcode}
StringLiteral str1 = "Some Phrase!";
StringLiteral str2 = str1.map(to_lower_case).pop_front<2>().front<7>();
// str1 equals "me phra"
\end{cppcode}
Если мы хотим использовать \verb"StringLiteral" в аргументах шаблона, то мы в частности хотим подставлять их в \verb"Value_t" в качестве первого аргумента.
Но тогда полезно завести псевдоним аналогично \verb"Int_t" для типа \verb"int".
А для поддержки операций еще и соответствующую переменную.
Будет это выглядеть так:
\begin{cppcode}
template<StringLiteral str>
using Str_t = Value_t<StringLiteral<str.capacity()>, str>;

template<StringLiteral str>
constexpr Str_t<str> str_v = {};
\end{cppcode}
Теперь у нас есть класс-wrapper и для строковых литералов, который будет совместим со списками типов.

\subsubsection{Операции}

\paragraph{Конкатенация строк}

Полезно переопределить оператор \verb"operator+" для конкатенации строк.
Для этого нужны три оператора:
\begin{cppcode}
template<std::size_t N, std::size_t M>
constexpr StringLiteral<N + M - 1> operator+(const StringLiteral<N>& left,
                                             const StringLiteral<M>& right) {
  char buffer[N + M - 1]{};
  std::copy_n(left.data, N - 1, buffer);
  std::copy_n(right.data, M, buffer + (N - 1));
  return buffer;
}

template<std::size_t N, std::size_t M>
constexpr StringLiteral<N + M - 1> operator+(const StringLiteral<N>& left,
                                             const char (&right)[M]) {
  return left + StringLiteral(right);
}

template<std::size_t N, std::size_t M>
constexpr StringLiteral<N + M - 1> operator+(const char (&left)[N],
                                             const StringLiteral<M>& right) {
  return StringLiteral(left) + right;
}
\end{cppcode}
Теперь можно делать такие вещи:
\begin{cppcode}
StringLiteral str1 = "Some";
StringLiteral str2 = "Phrase";
StringLiteral str3 = str1 + " " + str2 + "!";
// str3 equals "Some Phrase!"
\end{cppcode}
А если вспомнить про псевдонимы \verb"Str_t" и \verb"str_t", то можно сделать еще так:
\begin{cppcode}
auto str = str_v<"Some "> + str_v<"Phrase!">;
// str equals str_v<"Some Phrase!">
\end{cppcode}

\paragraph{Вывод в стандартный поток}

Так же полезно иметь возможность вывести строковый литерал на печать в стандартный поток вывода.
Для этого определим оператор:
\begin{cppcode}
template<std::size_t N>
std::ostream& operator<<(std::ostream& out, StringLiteral<N> str) {
  out << str.data;
  return out;
}
\end{cppcode}
После чего можно спокойно писать
\begin{cppcode}
StringLiteral str1 = "Some";
StringLiteral str2 = "Phrase";
StringLiteral str3 = str1 + " " + str2 + "!";
std::cout << str3 << '\n';  // prints "Some Phrase!"
\end{cppcode}

\subsubsection{Преобразование встроенных типов к StringLiteral}

\paragraph{Печать списка элементов в стандартный поток}

В начале давайте начну с того, для чего это нужно.
Пусть у нас есть список элементов, скажем
\begin{cppcode}
template<auto... Es>
struct List {};

using List1 = List<1, 3., 'c', true>;
\end{cppcode}
И мы хотим этот список распечатать в стандартный поток, то нужно составить соответствующую строку, которая бы представляла состояние списка.
Это можно сделать следующим образом:
\begin{enumerate}
\item Заменить каждый элемент списка на StringLiteral, чтобы потом оперировать только со строками.
Тогда должен получиться список вида
\begin{cppcode}
List<StringLiteral("1"), StringLiteral("3"), StringLiteral("c"), StringLiteral("true")>;
\end{cppcode}

\item Теперь эти строки надо склеить с помощью соединения \verb$", "$ и получить строку
\begin{cppcode}
"1, 3, c, true"
\end{cppcode}

\item После чего для красоты добавить открывающиеся и закрывающиеся скобки:
\begin{cppcode}
"[1, 3, c, true]"
\end{cppcode}
\end{enumerate}

\paragraph{Имплементация}

Давайте напишем оператор для печати списка из произвольных элементов
\begin{cppcode}
template<auto... Es>
struct List {};

template<class TList, class = std::enable_if_t<isList<TList>>>
std::ostream& operator<<(std::ostream& out, TList) {
  if constexpr (isEmpty<TList>) {
    out << "[ ]";
  } else {
    using SList = MapT<ToStringLiteral_t, TList>;
    static constexpr StringLiteral content =
        FoldL0T<JoinWithComma, SList>;
    static constexpr StringLiteral txt = "[" + content + "]";
    out << txt;
  }
  return out;
}
\end{cppcode}
При этом вспомогательные классы и методы выглядят так
\begin{cppcode}
template<StringLiteral left, StringLiteral right, StringLiteral delimiter>
struct Join {
  static constexpr auto value = left + delimiter + right;
};

template<StringLiteral left, StringLiteral right>
struct JoinWithComma : Join<left, right, ", "> {};
\end{cppcode}
Как это используется:
\begin{cppcode}
template<auto... Es>
struct List {};

using List1 = List<1, 3., 'c', true>;

std::cout << List1{} << '\n';
\end{cppcode}
То есть мы должны создать фиктивный объект типа \verb"List1" и отдать его на печать.

Теперь давайте пройдемся по коду оператора \verb"operator<<".
Посмотрим только на интересную часть печати не пустого списка:
\begin{enumerate}
\item В начале мы переводим каждый элемент из списка в строковый литерал, поэлементным применением:
\begin{cppcode}
using SList = MapT<ToStringLiteral_t, TList>;
\end{cppcode}

\item После чего объединяем все литералы в одну строку, соединяя элементы с помощью \verb$", "$:
\begin{cppcode}
static constexpr StringLiteral content = FoldL0T<JoinWithComma, SList>;
\end{cppcode}

\item Ну и наконец-то добавляем скобки по краям:
\begin{cppcode}
static constexpr StringLiteral txt = "[" + content + "]";
\end{cppcode}
\end{enumerate}

\paragraph{Проблемы}

И вот в этом всем мероприятии есть одна маленькая неприятная деталь: в C++ НЕТ \verb"constepxr" методов, которые бы переводили встроенные типы в строки.
Точнее есть метод \verb"to_chars", который умеет целочисленные типы переводить в строки.
И начиная с C++23 он является \verb"constexpr" только для целочисленных типов.
Если же вы хотите использовать его для \verb"float" или \verb"double", то эти преобразования вам придется писать руками.
Так себе занятие, учитывая возможные проблемы с потеряй точности.
Тем не менее, это сделать можно и этим мы сейчас займемся.
Будем страдать.

\paragraph{Вспомогательный тип \texttt{FixedString}}

Я хочу выделить отдельную \verb"constexpr" функцию, которая будет переводить значение встроенного типа в \verb"StringLiteral".
В идеале мы бы хотели что-то такое:
\begin{cppcode}
template<class T>
StringLiteral to_string(T val) {
  /*implementation*/
}
\end{cppcode}
Однако, \verb"StringLiteral" не является классом, это шаблон для которого надо указать параметр -- длина строки.
\begin{cppcode}
template<class T>
StringLiteral<?> to_string(T val) {
  /*implementation*/
}
\end{cppcode}
Так как это шаблонный параметр, то эта длина должна быть известна во время компиляции.
С другой стороны, длина строкового представления числа с плавающей запятой зависит от конкретного значения в переменной \verb"val".
И так как это аргумент функции, то конкретную длину мы можем узнать только в run-time.

Чтобы решить эту проблему нужно ввести еще один вспомогательный тип -- строку фиксированной длины \verb"FixedString".
Для каждого целого типа \verb"T" длина строкового представления будет уж точно не больше, чем длина бинарного представления числа.
То есть для \verb"int" должно быть достаточно $32$ знаков, а для \verb"long long" -- $64$ знака.
Если же дано число с плавающей запятой, то для \verb"float" точно можно оценить длину целой и дробной части в $32$ знака, а для \verb"double" в $64$ знака соответственно. 
Однако, у этого правила есть исключение.
Для типа \verb"char" мы все же хотим печатать не целое число, а соответствующий символ.
Аналогично для \verb"bool" мы хотим печатать не $0$ и $1$, а \verb"false" и \verb"true".

Вот как будет выглядеть вспомогательный тип:
\begin{cppcode}
template<size_t N>
struct FixedString {

  std::array<char, N> data{};
  size_t size = 0;

  constexpr FixedString() = default;

  template<size_t M, class = std::enable_if_t<(M <= N + 1)>>
  constexpr FixedString(const char (&buffer)[M]) {
    if (M == 0)
      return;
    // Not null terminated
    std::copy_n(buffer, M - 1, data.data());
    size = M - 1;
  }
  template<size_t M, class = std::enable_if_t<(M <= N)>>
  constexpr FixedString(char (&buffer)[M]) {
    std::copy_n(buffer, M, data.data());
    size = M;
  }

  constexpr std::string_view view() const {
    return {data.data(), size};
  }
  constexpr char back() const {
    return data[size - 1];
  }
  constexpr bool pop_back() {
    if (size == 0)
      return false;
    --size;
    return true;
  }
  constexpr bool append(char c) {
    if (size < N) {
      data[size++] = c;
      return true;
    }
    return false;
  };
  template<size_t M>
  constexpr bool append(const char (&buffer)[M]) {
    if (M == 0)
      return true;
    for (size_t i = 0; i < M - 1; ++i) {
      bool result = append(buffer[i]);
      if (!result)
        return false;
    }
    return true;
  }
  template<size_t M>
  constexpr bool append(const FixedString<M>& str) {
    for (size_t i = 0; i < str.size; ++i) {
      bool result = append(str.data[i]);
      if (!result)
        return false;
    }
    return true;
  }
};
\end{cppcode}
Данный тип хранит строку, которая не заканчивается \verb"'\0'".
Так как мы будем использовать эту строку как буффер, в который будут складываться символы, то делаем только методы добавления и удаления в конец, доступ с конца.
Если добавляемый элемент уже не влезает, мы его просто игнорируем.
Чтобы понять добавили мы элемент или нет, будем возвращать \verb"true" в случае, если добавили и \verb"false", если нет.
Да, я понимаю, что лучше \verb"enum class" со статусом, но мне было лень тут разводить это все.

\paragraph{Преобразование целых типов в StringLiteral}

Теперь как будет выглядеть процесс преобразования значения в строку.
Прежде всего у нас будет \verb"constexpr" функция \verb"to_string", которая будет переводить встроенные значения в литералы.
Вот так будет выглядеть случай целого типа отличного от \verb"bool" или \verb"char".
\begin{cppcode}
template<class T,
         class = std::enable_if_t<(std::is_integral_v<T> &&
                                   !isBool<T> && !isChar<T>)>>
constexpr FixedString<sizeof(T) * 8> to_string(T val) {
  FixedString<sizeof(T) * 8> str{};
  auto result = std::to_chars(&*str.data.begin(), &*str.data.end(), val);
  str.size = result.ptr - &*str.data.begin();
  return str;
}
\end{cppcode}
В этой имплементации я пользуюсь тем, что в C++23 функция \verb"std::to_char" является \verb"constexpr" для целых типов.
Чтобы обслужить случай \verb"bool" и \verb"char" мы добавим перегруженные версии функций так:
\begin{cppcode}
constexpr FixedString<8> to_string(bool val) {
  if (val)
    return "true";
  return "false";
}

constexpr FixedString<2> to_string(char val) {
  FixedString<2> str;
  str.append(val);
  return str;
}
\end{cppcode}
Самый интерес представляет данная функция для типов с плавающей запятой.
Это мы обсудим ниже, а сейчас я хочу собрать общее решение.
Чтобы суметь преобразовать \verb"FixedString" в \verb"StringLiteral", мы должны получать \verb"FixedString" в качестве шаблонного параметра, чтобы можно было считать его длину во время компиляции.
Для этого создадим вспомогательный класс \verb"Context" так:
\begin{cppcode}
template<FixedString str>
struct Context {
  static constexpr StringLiteral<str.size + 1> make() {
    char tmp[str.size + 1]{};
    std::copy_n(str.data.data(), str.size, tmp);
    tmp[str.size] = '\0';
    return tmp;
  }
};
\end{cppcode}
Теперь конвертация \verb"FixedString" в \verb"StringLiteral" можно сделать так
\begin{cppcode}
constepxr FixedString<64> str1 = "abc";
StringLiteral str2 = Context<str1>::make();
\end{cppcode}
Теперь сделаем удобную \verb"constexpr" функцию вместо класса:
\begin{cppcode}
template<auto val, class = enableIfIntegralOrFloat<val>>
constexpr auto toStringLiteral() {
  return Context<to_string(val)>::make();
}
\end{cppcode}
При этом мы включаем эту функцию только для целых встроенных типов и для встроенных типов с плавающей запятой.
\begin{cppcode}
template<auto val>
constexpr bool isIntegralOrFloat = std::is_integral_v<decltype(val)> ||
                                   std::is_floating_point_v<decltype(val)>;
template<auto val>
using enableIfIntegralOrFloat = std::enable_if_t<isIntegralOrFloat<val>>;
\end{cppcode}
Теперь сделаем класс-wrapper для хранения сконвертированного значения:
\begin{cppcode}
template<auto val>
struct ToStringLiteral_t
    : Value_t<decltype(toStringLiteral<val>()), toStringLiteral<val>()> {};
\end{cppcode}
И теперь осталось сделать еще $3$ варианта класса как раньше:
\begin{cppcode}
template<auto val>
constexpr auto ToStringLiteral = ToStringLiteral_t<val>::value;

template<class W>
struct ToStringLiteral_w : ToStringLiteral_t<W::value> {};

template<class W>
constexpr auto ToStringLiteral_u = ToStringLiteral_w<W>::value;
\end{cppcode}
Теперь если надо сконвертировать все элементы контейнера в StringLiteral, это можно сделать так:
\begin{cppcode}
template<auto... Es>
struct List {};

using List1 = List<1, 2., 'c', true>;
using List2 = Map<ToStringLiteral_t, List1>;
// List2 equals 
// List<StringLiteral("1"), StringLiteral("2"), StringLiteral("c"), StringLiteral("true")>
\end{cppcode}
Теперь осталось рассказать про имплементацию метода \verb"to_string" для типов с плавающей запятой.

\paragraph{Обработка \texttt{float} и \texttt{double}}

Начнем с самой функции:
\begin{cppcode}
template<class T, class = std::enable_if_t<std::is_floating_point_v<T>>>
constexpr auto to_string(T Value, int precision = PrecisionLimit<T>) {
  constexpr size_t str_size = SizeOfUIntFor<T> * 2;

  // not constexpr for MSVC
  if (std::isnan(Value)) {
    return FixedString<str_size>("nan");
  }
  // not constexpr for MSVC
  if (std::isinf(Value)) {
    if (Value < 0)
      return FixedString<str_size>("-inf");
    return FixedString<str_size>("inf");
  }

  precision = std::min<int>(precision, PrecisionLimit<T>);
  FixedString<str_size> res;
  T val = (Value < 0) ? (res.append('-'), -Value) : Value;
  res.append(i_part(val));
  res.append('.');
  res.append(f_part(val, precision));
  return trim(res);
}
\end{cppcode}
Давайте пройдемся по функции.
\begin{enumerate}
\item В начале обрабатываем особые случаи:
\begin{itemize}
\item Случай \verb"nan"

\item Случай \verb"inf"

\item Случай \verb"-inf"
\end{itemize}
Эти случаи можно проверить с помощью STL методов \verb"isnan" и \verb"isinf".
Однако, они не для всех компиляторов являются \verb"constexpr".
Потому их придется написать самостоятельно (видимо оставляю это как упражнение для пытливого читателя).

\item Дальше идет обработка самого числа.
В начале мы обрезаем точность, которую нам передали в качестве аргумента до максимальной точности данного типа.

\item Заводим \verb"FixedString" нужного размера (удвоенная битовая длина числа).

\item Отдельно записываем знак, если нужно.

\item Вычисляем целую часть и пишем ее в \verb"FixedString".

\item Добавляем \verb$"."$ после целой части.

\item Вычисляем дробную часть и пишем ее в \verb"FixedString" после точки.

\item Удаляем нули в конце дробной части.
Если дробная часть вся удалилась, то удаляем и знак точки.
\end{enumerate}
Теперь осталось понять, как вычислять целую и дробную часть.
Для целой части все просто:
\begin{cppcode}
template<class T, class = std::enable_if_t<std::is_floating_point_v<T>>>
constexpr FixedString<SizeOfUIntFor<T>> i_part(T value) {
  FixedString<SizeOfUIntFor<T>> res{};
  T val = (value < 0) ? (res.append('-'), -value) : value;
  UIntOf<T> ipart = static_cast<UIntOf<T>>(val);
  res.append(to_string(ipart));
  return res;
}
\end{cppcode}
Здесь \verb"T" либо \verb"float" либо \verb"double".
Сама функция работает так:
\begin{enumerate}
\item Заводим \verb"FixedString" размера равного количеству знаков в бинарном представлении числа.
Этого точно хватит.

\item Учитываем знак, если надо.

\item Теперь кастуем тип \verb"T" в подходящий целочисленный беззнаковый тип \verb"UIntOf<T>", что отрезает дробную часть.

\item И теперь добавляем в \verb"FixedString" конвертацию целой части, которая у нас уже имеется.
И возвращаем результат.
\end{enumerate}
Теперь дробная часть.
Здесь надо учитывать точность.
Для разных типов \verb"float" и \verb"double" это будут свои константы, которые можно настроить через шаблоны.
Я пропущу эти вспомогательные вещи.
Вот как выглядит функция:
\begin{cppcode}
template<class T, class = std::enable_if_t<std::is_floating_point_v<T>>>
constexpr FixedString<FractionSize<T>> f_part(T value, int precision) {
  T fpart = value - static_cast<UIntOf<T>>(value);
  precision = std::min<int>(precision, PrecisionLimit<T>);
  UIntOf<T> scaled = static_cast<UIntOf<T>>(fpart * pow10_table[precision] + 0.5);
  
  FixedString<FractionSize<T>> tmp{};
  tmp.size = precision;
  for (int i = precision - 1; i >= 0; --i) {
    tmp.data[i] = (scaled % 10) + '0';
    scaled /= 10;
  }
  return tmp;
}
\end{cppcode}
Функция работает следующим образом:
\begin{enumerate}
\item В начале вычленяем дробную часть \verb"fpart".

\item После этого умножаем дробную часть на $10^K$, где $K$ свое для каждого типа.
Для такого умножения я храню таблицу степеней $10$.
\begin{cppcode}[numbers = none]
inline constexpr std::array<double, 17> pow10_table = {
    1e0, 1e1,  1e2,  1e3,  1e4,  1e5,  1e6,  1e7, 1e8,
    1e9, 1e10, 1e11, 1e12, 1e13, 1e14, 1e15, 1e16};
\end{cppcode}

\item Преобразованием к целому типу я отрезаю первые \verb"precision" знаков у дробной части в виде целого числа.

\item Далее посимвольно записываю дробную часть в \verb"FixedString" и возвращаю результат.
\end{enumerate}
Для типа \verb"double" значение \verb"precision" берется равным $14$, а для типа \verb"float" равным $5$.

И осталось привести имплементацию функции \verb"trim", отрезающей нули в конце строкового представления числа с плавающей запятой.
\begin{cppcode}
template<size_t M>
constexpr auto trim(FixedString<M> str) {
  while (str.size > 0 && str.back() == '0')
    str.pop_back();
  if (str.size > 0 && str.back() == '.')
    str.pop_back();
  return str;
}
\end{cppcode}
В целом тут можно не удалять точку, если хочется подчеркнуть, что число с плавающей запятой.
